%% kapa: exam=1 round=24 year=2013 period=2 form=A subject=accounting subject_ko=회계학 seq=32 number=32
%% source: Q-Net 감정평가사 기출문제 (cst00309) artlSeq=5073802 file=extracted/1008/5073802/2013년 제24회 감정평가사 1차시험 기출문제/2교시(A형).hwp
%% conversion: kordoc 4.24.51 -> tools/gen1.py (auto); review=proofread-sonnet; flags=none
\begin{question}{32}
감평회계법인은 컨설팅과 회계감사서비스를 제공하고 있다. 지금까지 감평회계법인은 일반관리비 ￦270,000을 용역제공시간을 기준으로 컨설팅과 회계감사서비스에 각각 45\%와 55\%씩 배부해 왔다. 앞으로 감평회계법인이 활동기준원가계산을 적용하기 위해, 활동별로 일반관리비와 원가동인을 파악한 결과는 다음과 같다. \par
{\small\centering
\begin{tabular}{|c|c|c|}
\hline
\makecell{활동} & \makecell{일반관리비} & \makecell{원가동인} \\
\hline
\makecell{스탭지원} & \makecell{￦200,000} & \makecell{스탭수} \\
\hline
\makecell{컴퓨터지원} & \makecell{50,000} & \makecell{컴퓨터사용시간} \\
\hline
\makecell{고객지원} & \makecell{20,000} & \makecell{고객수} \\
\hline
\makecell{합계} & \makecell{￦270,000} &  \\
\hline
\end{tabular}\par}

컨설팅은 스탭수 35\%, 컴퓨터사용시간 30\% 그리고 고객수 20\%를 소비하고 있다. 활동기준원가계산을 이용하여 컨설팅에 집계한 일반관리비는 이전 방법을 사용하는 경우보다 얼마만큼 증가 또는 감소하는가? \par
\begin{choices}
  \item ￦32,500 감소
  \item ￦32,500 증가
  \item ￦59,500 감소
  \item ￦59,500 증가
  \item 변화 없음
\end{choices}
\end{question}
\answer{1}
